\section{Two conductor-eight quartic class certificates at the ramified cubic prime}
\label{sec:quadratic8-quartic-class-certificates}

\paragraph{Scope and fixed data.}
Put $p=269$, $K=\mathbb Q(\omega)$, and $t=\sqrt[3]p>0$.
Let $\psi=\omega_p^{67}$ be the primitive \emph{odd} quartic
character of conductor $p$, with $\psi(2)=-i$ and
$i\equiv82\pmod p$.
For $s=1,-1$ let $\delta_s=\chi_{8s}$ be the primitive quadratic
character of conductor eight, and set $\xi_s=\delta_s\psi$.
The signs in this paragraph label two different cyclic quartic
fields, not the inverse characters of a single field.
The inverse of $\xi_s$ cuts out the same field as $\xi_s$.

The exact unit coordinates, monic relations, residue maps and
interval outputs named below are part of this finite certificate.
The preliminary PARI class data are used only to propose the two
integers tested by rational isolation.
No complete fundamental-unit basis, conditional class computation
or new local strict-vanishing theorem is assumed.

\begin{theorem}[AC-G158-class: the two conductor-eight quartic class certificates]
\label{AC-G158-class}
For $s=1,-1$ define
\[
 \begin{gathered}
 F_s=\mathbb Q(w_s),\qquad
 w_s^4+2spw_s^2+100p=0,\qquad
 L_s=F_s(t),\qquad M_s=KL_s.
 \end{gathered}
 \tag{Q8QC.1}
\]
Then $F_s$ is the cyclic quartic field of $\xi_s$, its quadratic
subfield is $\mathbb Q(\sqrt p)$, and
\[
 \begin{gathered}
 D_{F_s}=2^6p^3,\\
 |D_{L_s}|=2^{18}3^4p^{11}\\
 =11331829223512982224732649157820416.
 \end{gathered}
 \tag{Q8QC.2}
\]
Both $L_s$ have degree twelve and only the roots of unity $\{\pm1\}$.
Their signatures, unit ranks and certified relations are
\[
 \begin{array}{c|c|c|c}
 s&(r_1,r_2)&r&h_{L_s}\\ \hline
 1&(0,6)&5&574m_1\\
 -1&(4,4)&7&28m_{-1},
 \end{array}
 \qquad p\nmid m_1m_{-1}.
 \tag{Q8QC.3}
\]
Here $m_s$ is the index, modulo torsion, of the explicitly
certified units in the full unit lattice.
In particular
\[
             \operatorname{Cl}(L_s)[p]=
             \operatorname{Cl}(F_s)[p]=0.
 \tag{Q8QC.4}
\]
The nontrivial $C_3$ part of $\operatorname{Cl}(M_s)[p]$ is zero.
No assertion about the remaining cubic-trivial part of this
normal-closure class group is made.
\end{theorem}

\begin{proof}
\proofstep{1. Exact field and character dictionary.}
The previously certified Gaussian-period calculation for $p=269$
provides the odd quartic field
\[
 \begin{gathered}
 F_0=\mathbb Q(v),\qquad v^4+pv^2+25p=0,\\
 v^2=-\frac{p+13\sqrt p}{2},\qquad
 g(v)=-\frac{v^3+219v}{65},\qquad
 g^2(v)=-v,\qquad\psi(g)=-i .
 \end{gathered}
 \tag{Q8QC.5}
\]
We retain the exact period enumeration, polynomial multiplication
and orientation from AC-G147-class and AC-G156-class.
This is a finite source input, not a new appeal to a numerical
period recognition.

The unique quadratic subfield of $F_0$ is $\mathbb Q(\sqrt p)$,
so it is disjoint from $\mathbb Q(\sqrt{2s})$.
In their degree-eight compositum put
\[
                         w_s=\sqrt{2s}\,v .
 \tag{Q8QC.6}
\]
The single ordered pair consisting of the nontrivial quadratic
automorphism and $g^2$ fixes $w_s$.
Its order-two subgroup is exactly the kernel of
$\chi_{8s}\psi$ on $C_2\times C_4$.
The polynomial (Q8QC.1) follows from (Q8QC.5);
it is Eisenstein at $p$ and has degree four.
Consequently its root generates exactly the desired fixed field.
Both $\xi_s$ are primitive of conductor $8p$: the factor at two
is primitive of conductor eight and the factor at $p$ has order
four, so neither local factor can cancel.
Moreover $\xi_s^2=\tau_p$, the quadratic character of conductor $p$.
The conductor-discriminant formula gives
$D_{F_s}=p(8p)^2$.

The character $\delta_1$ is even and $\delta_{-1}$ is odd.
Thus $\xi_1$ is odd and $F_1$ is imaginary, whereas $\xi_{-1}$
is even and $F_{-1}$ is totally real.
Their roots have squared absolute values
$p+13\sqrt p$ and $p-13\sqrt p$; both are positive since
$269>169$.
The automorphism $b$ fixing $\sqrt{2s}$ in the compositum has
\[
 b(w_s)=-\frac{s\,w_s^3}{130}-\frac{219w_s}{65},
 \qquad b^2(w_s)=-w_s,\qquad \xi_s(b)=-i .
 \tag{Q8QC.7}
\]
Equivalently it is represented on the character group by the
CRT class $809$, which is one modulo eight and two modulo $269$.
The exact verifier checks both polynomial identities, $b^2=-1$,
and the displayed character orientation.

The normal closure of $\mathbb Q(t)$ has group $S_3$ and unique
quadratic subfield $K$.
The cyclic quartic $F_s$ has the different quadratic subfield
$\mathbb Q(\sqrt p)$.
Their intersection is therefore $\mathbb Q$, giving
\[
 \operatorname{Gal}(M_s/\mathbb Q)=S_3\times C_4,\qquad
 [L_s:\mathbb Q]=12,\qquad[M_s:\mathbb Q]=24.
 \tag{Q8QC.8}
\]
For $s=1$ there are no real embeddings of $L_s$.
For $s=-1$ its four real embeddings choose the real cubic root
and any of the four real quartic roots; the other eight
embeddings form four complex pairs.
This proves the signatures in (Q8QC.3).

At $p$ the cubic and quartic completions are totally ramified of
coprime degrees three and four.  Their compositum is totally
ramified of degree twelve and is tame, so its discriminant
exponent is eleven.
At two the pure cubic algebra is unramified of total degree
three, since its polynomial discriminant is odd.
Unramified base change in that degree-three algebra multiplies
the quartic discriminant exponent six by three.
At three, $269$ is a cube in $\mathbb Q_3$: for
$f(X)=X^3-269$ one has $v_3(f(-1))=3>2v_3(f'(-1))=2$.
The cubic algebra is therefore
$\mathbb Q_3\times\mathbb Q_3(\sqrt{-3})$, with discriminant
exponent one.  The quartic algebra is unramified at three,
so the exponent in $L_s$ is four.
There is no ramification elsewhere.  This proves (Q8QC.2).

The maximal abelian subfield of $L_s/\mathbb Q$ is $F_s$:
in (Q8QC.8) the normal subgroup generated by a reflection
fixing $L_s$ is the whole $S_3$ factor.
Every root of unity in $L_s$ therefore belongs to $F_s$.
An odd prime divisor of its order could only be $p$, since
$F_s$ is unramified outside $2,p$, but
$[\mathbb Q(\mu_p):\mathbb Q]=268>4$.
A root of unity of two-power order greater than two would
contain $\mathbb Q(i)$, incompatible with the unique real
quadratic subfield of $F_s$.
Thus the torsion group is exactly $\{\pm1\}$.

\proofstep{2. Exact units and the integral lattice indices.}
For each sign the candidate extraction records the raw generator
$x=w_s-t$, its degree-twelve polynomial, and the exact inverse
maps for $t,w_s$.
The extraction uses \texttt{bnfinit} only to suggest units and
a class-number selector.
No assertion from \texttt{bnfcertify}, or a conditional class
group produced by the extraction, is used.

The independently checkable algebra is the tensor basis
\[
 \{t^a w_s^b:0\leq a\leq2,\ 0\leq b\leq3\},
 \qquad\text{index }4a+b .
 \tag{Q8QC.9}
\]
The script \path{verify_units_exact.py} constructs the two
rational regular-representation matrices $T,W$ in this basis
using $T^3=p$ and $W^4+2spW^2+100p=0$.
It checks $TW=WT$, the characteristic polynomial of $W-T$,
its irreducibility and both raw inverse maps.
For every proposed unit it evaluates the full raw polynomial
at $W-T$, recovers its exact tensor coordinates, computes its
characteristic polynomial and separately verifies that polynomial
by matrix Horner evaluation.
Each polynomial is monic integral with constant coefficient one.
Thus each listed element and its inverse are integral.
All rational denominators are kept throughout these exact steps.

The complete unit coordinates and relations are the files
\path{chi8/all_units_exact_certificate.json} and
\path{chi_minus8/all_units_exact_certificate.json}.
Their SHA256 values are, respectively,
\begin{quote}\ttfamily\small
cc09226d6f5a1072782b152ac8161a701\\
47000943a6a75d7642e54d2ea889645,\\
d3c36fd4771a4cafca7fd0f730591735f\\
cb77990e7edd7a960de22b40f022601.
\end{quote}
The coordinate denominators, in their fixed unit order, are
\[
 \begin{aligned}
 s=1:&\quad13,\ 1,\ 3497,\ 52455,\ 52455,\\
 s=-1:&\quad13,\ 130,\ 130,\ 3497,\ 3497,\ 104910,\ 52455.
 \end{aligned}
 \tag{Q8QC.10}
\]
These data include all twelve units, not just a selected local
witness.  In particular the denominators divisible by $p$ in
(Q8QC.10) have not been suppressed.

To certify independence and $p$-saturation, use the rational
prime $\ell=206593$.
Complete trial division proves primality.
Each cubic and quartic polynomial has its full set of simple
roots modulo $\ell$, giving twelve degree-one residue maps for
each sign.
The exact script finds a primitive root $g$ modulo $\ell$,
verifying every prime-divisor test for $\ell-1$, and sets
$h=g^{(\ell-1)/p}$.
Then $h$ has exact order $p$.
For every unit image $u$ the equation
\[
                u^{(\ell-1)/p}=h^{\,a(u)}
 \tag{Q8QC.11}
\]
defines a homomorphism to $\mathbb F_p$.
It kills $p$th powers and the two-element torsion subgroup.
Simple-root Hensel lifting justifies these residue maps on the
full ring of integers; every displayed unit denominator is
invertible at $\ell$.

The first sign has the selected five by five integer minor
\[
 \begin{pmatrix}
199&85&90&159&233\\
70&85&264&233&110\\
70&85&264&36&159\\
199&85&90&110&36\\
199&214&153&235&181
 \end{pmatrix},
 \quad
 \det=-217327075290\equiv246\pmod{269}.
 \tag{Q8QC.12}
\]
For the second sign a selected seven by seven minor is
\[
 \begin{pmatrix}
199&0&56&264&90&63&210\\
70&213&0&90&264&124&63\\
70&56&0&90&264&210&226\\
199&0&213&264&90&226&124\\
199&0&56&61&153&248&9\\
70&213&0&153&61&122&248\\
70&56&0&153&61&9&104
 \end{pmatrix}.
 \tag{Q8QC.13a}
\]
\[
 \begin{gathered}
 \det=4197143958599830,\\
 \det\equiv142\pmod{269}.
 \end{gathered}
 \tag{Q8QC.13}
\]
The complete twelve-row matrices, unit values and discrete
logarithms are retained in each
\path{269_saturation_certificate.json}.

A primitive integer relation among the units would reduce to
a nonzero vector in the kernel of the corresponding matrix,
contradicting (Q8QC.12) or (Q8QC.13).
They are therefore independent.
Their number equals the full Dirichlet rank, so their subgroup
has finite index $m_s$ modulo torsion.
The same full-rank map modulo $p$ implies $p\nmid m_s$ by the
Smith normal form of this lattice inclusion.
This establishes precisely the index information used below;
it does not assert $m_s=1$.

\proofstep{3. Weighted regulator normalization.}
Let $\rho=p^{1/3}$ and put
\[
             A=\sqrt{p+13\sqrt p},\qquad
             B=\sqrt{p-13\sqrt p}.
 \tag{Q8QC.14}
\]
For $s=1$, use the six complex-place representatives
$(t,w)=(\rho\omega^j,iA),(\rho\omega^j,iB)$, $0\leq j<3$.
Every row has weight two.
For $s=-1$, use the four real embeddings
$(\rho,A),(\rho,-A),(\rho,B),(\rho,-B)$ with weight one,
and the four complex representatives with $\rho$ replaced by
$\rho\omega$, with weight two.

For a unit $U$ and embedding of weight $d_v$, the entry is
$d_v\log|\sigma_v(U)|$.
The rows sum to zero by the product formula.
Delete the last row and take the absolute determinant of the
remaining rank by rank matrix.
This is the projected weighted-log convention of the retained
analytic class-number formula, and gives
\[
                       R_{\mathcal U_s}=m_sR_{L_s}.
 \tag{Q8QC.15}
\]
The computation uses the full exact tensor expression at every
embedding, including all denominators.
The script asserts that each squared absolute value is positive,
that each column sum encloses zero and that the determinant
excludes zero.
Its regulator precisions are 2500 bits for $s=1$ and
1500 bits for $s=-1$; all subsequent analytic operations use
256-bit outward interval arithmetic.

\proofstep{4. The complete zeta factorization and Dirichlet factors.}
Let $\rho_{\mathrm{std}}$ be the two-dimensional standard Artin
representation of the $S_3$ field $K(t)$.
The permutation representation on the embeddings of $L_s$
is
$(1\oplus\rho_{\mathrm{std}})\otimes
 (1\oplus\tau_p\oplus\xi_s\oplus\bar\xi_s)$.
It follows at every Euler factor that
\[
 \begin{aligned}
 \zeta_{L_s}(z)&=\zeta(z)L(z,\tau_p)|L(z,\xi_s)|^2\\
       &\quad\cdot L(z,\rho_{\mathrm{std}})
       L(z,\rho_{\mathrm{std}}\tau_p)
       |L(z,\rho_{\mathrm{std}}\xi_s)|^2 ,
 \end{aligned}
 \tag{Q8QC.16}
\]
where the squared absolute notation at real $z$ denotes the
product with the conjugate-character factor.

Write $a_n$ for the standard Artin coefficients.
For a good prime $q$, their trace is the number of roots of
$X^3-p$ modulo $q$ minus one.
For $q\equiv2\pmod3$ the local coefficients are one in even
degree and zero in odd degree.
For $q\equiv1\pmod3$ they satisfy
$a_{q^r}=a_q a_{q^{r-1}}-a_{q^{r-2}}$, with $a_q=2$ for
a cubic residue and $a_q=-1$ otherwise.
At three the local coefficients are all one; at $p$ all
positive-degree local coefficients vanish.
These descriptions include the bad Euler factors: at three
there is one invariant line with trivial eigenvalue, and
at $p$ the cubic inertia has no invariant vector.
They give the multiplicative bound
\[
                              |a_n|\leq d(n)\leq n.
 \tag{Q8QC.17}
\]
Multiplication by each Dirichlet character supplies the exact
twisted coefficients, with zero value at its conductor primes.

For a primitive even character $\nu$ of conductor $f$, set
\[
       S_\nu=\sum_{a=1}^{f-1}\bar\nu(a)
                     \log\left(2\sin\frac{\pi a}{f}\right).
\]
The retained finite Fourier formula gives
$|L(1,\nu)|^2=|S_\nu|^2/f$.
For a primitive odd character the corresponding formula is
\[
 |L(1,\nu)|^2=\frac{\pi^2}{f}|B_{1,\nu}|^2,\qquad
 B_{1,\nu}=\frac1f\sum_{a=1}^{f-1}a\nu(a).
 \tag{Q8QC.18}
\]
Here direct exact summation gives
\[
                       B_{1,\xi_1}=8+10i.
 \tag{Q8QC.19}
\]
For the even character $\xi_{-1}$, the script uses the finite
log-sine expression with $f=2152$.
Using $\nu$ rather than $\bar\nu$ conjugates its sum and leaves
the required squared norm unchanged.

For $\tau_p$, the positive quadratic Gauss sign is fixed by the
retained finite cosine-sum certificate and rechecked by the
outward interval evaluation.
The full value, including its sign, is therefore
\[
 L(1,\tau_p)=-\frac1{\sqrt p}
       \sum_{a=1}^{p-1}\tau_p(a)
                      \log\left(2\sin\frac{\pi a}{p}\right).
 \tag{Q8QC.20}
\]
The positive interval in the output is used directly.
No quadratic class-number or fundamental-unit assumption enters
this Dirichlet factor.

\proofstep{5. Primitive Hecke conductors, including the wild prime two.}
The standard representation is induced from the nontrivial cubic
character of $K(t)/K$.
Over $K$ its conductor is $p\mathcal O_K$: it is tamely nontrivial
at the inert prime $p$ and unramified elsewhere, including three
by the local cube argument above.
The class number of $K$ is one, since its imaginary quadratic
Minkowski bound is $2\sqrt3/\pi<2$.
Thus each finite Hecke character here can be expressed as a
character of the appropriate residue ring, trivial on the six
units of $\mathcal O_K$.

The three Hecke twists corresponding to
$\rho_{\mathrm{std}}$, $\rho_{\mathrm{std}}\tau_p$ and
$\rho_{\mathrm{std}}\xi_s$ have moduli
\[
              p\mathcal O_K,\quad p\mathcal O_K,\quad
              8p\mathcal O_K,
 \qquad
 Q=3p^2,\ 3p^2,\ 3(8p)^2 .
 \tag{Q8QC.21}
\]
At $p$, the residual characters have orders three, six and twelve:
the norm from $\mathbb F_{p^2}^{\times}$ onto
$\mathbb F_p^\times$ is surjective, so the orders two and four
of the rational twists are retained and cannot cancel order three.

At two, $K_2/\mathbb Q_2$ is unramified quadratic.
For $a\in\mathcal O_{K_2}$,
\[
             \operatorname{Norm}(1+4a)
                         \equiv1+4\operatorname{Tr}(a)\pmod8.
 \tag{Q8QC.22}
\]
The residue trace $\mathbb F_4\to\mathbb F_2$ is surjective.
Choose a trace-one residue; its norm is five modulo eight,
and both primitive characters $\chi_8,\chi_{-8}$ take value
minus one on five.
The norm-pulled-back character is therefore nontrivial on
$1+4\mathcal O_{K_2}$ and trivial on $1+8\mathcal O_{K_2}$.
This proves the exact exponent three, rather than merely an
upper bound on the conductor.

Here is the finite Fourier step for these actual moduli.
Write $\mathcal O=\mathcal O_K$ and $f=p$ or $8p$.
The inverse different is $(1/\sqrt{-3})\mathcal O$.
The additive pairing
\[
 (r,\beta)\longmapsto
       \exp\left(\frac{2\pi i}{f}
          \operatorname{Tr}_{K/\mathbb Q}
                        \frac{-r\bar\beta}{\sqrt{-3}}\right)
 \tag{Q8QC.23}
\]
is perfect on $\mathcal O/f\mathcal O$.
In particular it remains perfect at two: in the Euclidean dual,
the factor two in $2/\sqrt{-3}$ cancels the factor one-half in
$\operatorname{Re}(z)=\operatorname{Tr}(z)/2$.
No assertion that two is a unit modulo eight is needed.

Extend a primitive multiplicative character $\theta$ by zero
on nonunits and let $G_\theta(\beta)$ be its finite Fourier sum.
For a nonunit frequency at $p$, choose a residue unit on which
the primitive character is nontrivial and use multiplication
by that unit to force the sum to zero.
For a frequency divisible by two in the modulus-eight factor,
choose $u=1+4a$ as in (Q8QC.22).
Then $(u-1)\beta=0$ modulo eight but $\theta(u)=-1$.
CRT makes the auxiliary unit one at the other prime.
The same multiplication argument gives zero.
For a unit frequency, change variables to obtain
\[
       G_\theta(\beta)=\theta(\bar\beta)^{-1}G_\theta(1).
 \tag{Q8QC.24}
\]
Parseval for the perfect finite pairing, together with the zero
nonunit frequencies, gives $|G_\theta(1)|=f$.
This proves the primitive Gauss step also at the prime power
$2^3$, without a squarefree-modulus shortcut.

The covolume of $f\mathcal O$ is $f^2\sqrt3/2$ and its
Euclidean dual is $(2/(f\sqrt{-3}))\mathcal O$.
For $a=2\pi/(\sqrt3 f)$ the Gaussian Fourier transform therefore
has scalar $1/(fy)$ and dual exponent $-a|\beta|^2/y$.
In (Q8QC.24), conjugating the dual variable changes
$\theta(\bar\beta)^{-1}$ to $\bar\theta(\beta)$.
Poisson summation, with the common six-unit factor on both
sides, yields
\[
 \Theta_\theta(y)=\epsilon y^{-1}
                         \Theta_{\bar\theta}(1/y),
 \qquad |\epsilon|=1,\qquad
 \Theta_\theta(y)=\sum_{n\geq1}a_n\nu(n)e^{-any}.
 \tag{Q8QC.25}
\]
The zero-frequency term is zero, since the character is nontrivial.
This establishes the precise normalization used in the interval
calculation.

\proofstep{6. Finite intervals and rational class-number isolation.}
Put $T=\Theta_\theta(1)$.
The computed interval excludes zero in all cases.
At $y=1$, (Q8QC.25) gives
$\epsilon=T/\bar T$.
For a real coefficient series this is exactly one.
Splitting the Mellin integral at one gives
\[
 \begin{aligned}
 L(1,\rho_{\mathrm{std}}\nu)
     &=A+\epsilon a\bar B,\\
 A&=\sum_{n\geq1}\frac{a_n\nu(n)e^{-an}}n,\qquad
 B=\sum_{n\geq1}a_n\nu(n)E_1(an).
 \end{aligned}
 \tag{Q8QC.26}
\]
Thus no root-number phase is guessed.
If the sums are stopped at $N$ and $z=e^{-a}$, (Q8QC.17) gives
the absolute tail bounds
\[
 \frac{e^{-a(N+1)}}{1-z},\qquad
 \frac{e^{-a(N+1)}}{a(1-z)},\qquad
 \frac{e^{-a(N+1)}((N+1)-Nz)}{(1-z)^2}
 \tag{Q8QC.27}
\]
for $A,B,T$, respectively.
The middle inequality uses $E_1(x)\leq e^{-x}/x$.
These radii are added to every relevant real or imaginary
component before any division or multiplication.

The single analytic script is
\path{certify_quadratic8_quartic_analytic.py}.
It takes a case directory and the explicit label
\texttt{chi8} or \texttt{chi\_minus8}.
It uses $N=12000$ for the first two Hecke factors and $N=96000$
for the conductor-$8p$ factor.
The exact arithmetic coefficient sieve, finite Dirichlet sums,
tail radii, nonzero theta checks, phase divisions and final
intervals are saved separately for each case.

For completeness, the denominator bound used in the isolation
does not depend on a regulator lower bound.
For degree $n=12$ and signature $(r_1,r_2)$ put
\[
 \begin{gathered}
 U=\left\lceil
 \frac{12!}{12^{12}}\left(\frac43\right)^{r_2}
                     (\lfloor\sqrt D\rfloor+1)\right\rceil,\\
 H=4^{12}\left\lceil U^{4/3}\right\rceil .
 \end{gathered}
 \tag{Q8QC.28}
\]
Minkowski gives an integral ideal representative of norm at most
$U$ in each class, since $\pi>3$.
The count of such ideals is at most
$U^{4/3}\zeta_{L_s}(4/3)$.
Each rational Euler factor is bounded by the twelfth power of
the Riemann zeta factor, and
$\zeta(4/3)\leq1+\int_1^\infty x^{-4/3}\,dx=4$.
Hence $h_{L_s}\leq H$.
The script computes the ceiling in $H$ by an exact integer
cube-root search, with adjacent-integer inequalities checked.

In the present cases the exact integers are
\[
 \begin{aligned}
 U_1&=32132494478729,\\
 H_1&=17138750297576425136324608,\\
 U_{-1}&=18074528144285,\\
 H_{-1}&=7958096243376782202372096.
 \end{aligned}
 \tag{Q8QC.29}
\]
The retained analytic class-number formula, with two roots of
unity, reads
\[
 h_{L_s}R_{L_s}
   =\frac{2\sqrt D}{2^{r_1}(2\pi)^{r_2}}
                      \operatorname*{Res}_{z=1}\zeta_{L_s}(z).
 \tag{Q8QC.30}
\]
Inserting (Q8QC.16), the weighted regulator (Q8QC.15), and all
outward intervals yields
\[
 \begin{array}{c|c|c|c}
 s&c_s\ (r_s>1)&|r_s-c_s|&1/H_s\\ \hline
 1&574&<1.30\cdot10^{-56}&>5.83\cdot10^{-26}\\
 -1&28&<6.65\cdot10^{-58}&>1.25\cdot10^{-25}.
 \end{array}
 \tag{Q8QC.31}
\]
Here $r_s=h_{L_s}/m_s$.  The second column specifies the integer
selector, not a claim
that $r_s$ has been assumed equal to it.

Because $r_s>1$, the positive integer index satisfies
$m_s<h_{L_s}\leq H_s$.
If $h_{L_s}\ne c_sm_s$, then
$|r_s-c_s|\geq1/m_s>1/H_s$, contradicting (Q8QC.31).
Thus $h_{L_1}=574m_1$ and $h_{L_{-1}}=28m_{-1}$.
Neither selector is divisible by $269$, and the same is true
of the indices by Step2.
This proves the two $L_s$ class assertions in (Q8QC.4).
The computation proves these relations and prime-to-$p$ statements;
it does not prove that the class numbers are exactly $574$ or $28$.

\proofstep{7. The precise descended class sectors.}
The extension $L_s/F_s$ has degree three.
Extension of ideal classes followed by norm is multiplication
by three, so the vanishing for $L_s$ implies
$\operatorname{Cl}(F_s)[p]=0$.

In $M_s/\mathbb Q$ let $r$ be the reflection fixing $L_s$.
On $\operatorname{Cl}(M_s)[p]$, norm to $L_s$ followed by
extension is $1+r$, and is zero.
Every conjugate reflection therefore acts by minus one.
Products of reflections generate the cubic rotations, which
must act trivially.
This proves that the nontrivial $C_3$ part is zero.
It supplies no assertion about the remaining cubic-trivial part.
\end{proof}

\paragraph{Original coefficient boundary and source scope.}
For the actual finite labels $\chi^{-j}\chi_{\pm8}$ with $j=1,2$,
the residual Kummer characters at effective weights $66$ and
$200$ modulo $268$ are respectively
$\bar\chi^{-j}\bar\xi_s^{-1}$ and
$\bar\chi^{-j}\bar\xi_s$.
Their selected cubic sectors are exactly the sectors proved
zero above.  The quartic character at $269$ is a residual
splitting character, not an enlargement of the finite label
in the actual host.
A separate proof must still show local unit injectivity, lift
over the original possibly ramified coefficient ring, and pay
the actual common-window local-index cost.

The retained primary analytic normalizations are the analytic
class-number formula in the weighted-log convention
(MIT 18.785, 2019 Lecture19, pp.1 and4), and the finite
Dirichlet identities at DLMF25.15.3,5--6 at their previously
audited scopes.  The standard class-field correspondence,
conductor-discriminant formula, Minkowski bound, Dirichlet unit
theorem and Poisson summation retain their previously audited
source scopes.
The exact Gaussian-period dictionary is retained from the
complete earlier certificate.
Steps1--7 above give the new field, local conductor,
prime-power Fourier, regulator-index and rational-isolation
arguments in full.
Arb supplies the previously audited outward interval operations;
its package version, script hashes, precision and exact input
hashes are recorded in the outputs.
No fresh whole proof of these deep retained sources, complete
unit basis, full normal-closure class vanishing, strict theorem,
all-prime estimate, H04.4 conclusion or Lean execution is claimed.
